\section{Upper bounds for the theta lifts}\label{sec:upper-bound}

The ramified source subspace \(S\) provides forms whose lifts have
controlled length and whose images in \(V\) have controlled index.
The index estimate applies to all twists allowed at each finite place.
The length estimate concerns the trivial twist.

\begin{theorem}[Upper bounds]\label{thm:upper-bound}
\begin{enumerate}
\item The forms \(\phi\in S\) that are integral at every finite place and satisfy
\[
  \|\Theta_1\mathfrak p_g\phi\|_\sigma^2
  \le \exp\!\bigl(H_0-m\mathcal R+E_4\bigr)
\]
for all \(g\in\overline G\) and \(\sigma\in\operatorname{Aut}(\mathbb C)\)
span \(S\).
\item Let \(\mathcal T\subseteq\overline G\) be finite.
Suppose that every member \(y_i\) of a family can be written as
\(y_i=\mathfrak p_{g_i}\phi_i\), with \(g_i\in\mathcal T\) and
\(\phi_i\in S\) integral at every finite place.
Then the logarithmic index of \((y_i)_i\) is at most \(2|\mathcal T|(\dim S)E_3\).
\end{enumerate}
\end{theorem}

\subsection{The seesaw identity and the projection}

The principal-level forms on \(Y_0\) carry an action of \(\overline G\),
the product of the general linear groups over the residue fields at the
selected places above \(M\). The central idempotent \(e_0\) of its group
algebra projects onto \(V\). Write \(\rho\) for the group action and
\(\mathcal P\) for the Poincar\'e series map. The lower chain of lifts
from principal level, along \(Y_0\to\cdots\to Y_m\), is denoted by \(L_\xi\).
Its restriction to \(V\) is \(\Theta_\xi\).

\emph{Unfolding the integrals of the lower chain.}
Let \(F=E^+\), and take adeles over \(F\) here.
Write \(G_j=U(Y_j)\) and \([G_j]=G_j(F)\backslash G_j(\mathbb A)\). Let \(J^\circ_j\) be the stabilizer of the standard
lattice at the finite places and \(J_{\xi,j}\subset J^\circ_j\) the
level for the twist \(\xi\).
The lower theta kernel \(K_{\xi,j}(g,h)\) is the theta value of the
Schwartz datum after the Weil action of
\((g,h)\in G_j(\mathbb A)\times G_{j+1}(\mathbb A)\).
The lower lift \(L_{\xi,j}\) integrates \(K_{\xi,j}(g,h)f(g)\)
over \([G_j]\) and multiplies by \(d_{\xi,j}=[J^\circ_j:J_{\xi,j}]\).
For a form \(z\) of conjugate type, put
\[
 L_{\xi,j}^{\mathrm t}z(g)=d_{\xi,j}
          \int_{[G_{j+1}]}K_{\xi,j}(g,h)z(h)\,dh .
\]
Exchanging the order of integration gives
\[
  \int_{[G_{j+1}]} z\,L_{\xi,j} f
      = \int_{[G_j]}(L_{\xi,j}^{\mathrm t}z)f .
\]
Apply this identity from \(j=m-1\) down to \(0\).
The pairing with the test form \(x_m\) becomes the integral of
\(\mathcal P\phi\) against the successive descents of \(\overline{x_m}\).
Rapid decrease gives the absolute convergence used in this step.
\leansrc{New/UpperBound/Seesaw/SeesawIdentity}{Transpose of the lift and absolute convergence}.

Let \(R_j\) be the Bessel subgroup for the grading of \(Y_j\).
The initial function \(\mathcal W_0\) on \([R_0]\) extends the
matrix coefficient of \(\phi\) on \(X_0\) by the prescribed
unipotent character. After \(j\) unfoldings, write \(\mathcal W_j\)
for the resulting function with the earlier outer variables fixed.
Fix the next outer variable \(g_j=(a_j,y_j)\), its infinite and
finite components. Put \(W_{j,g_j}(r)=\mathcal W_j(rg_j)\) for \(j>0\)
and \(W_{0,g_0}(r)=\mathcal W_0(r)\) at the initial step.
Integrate this function times the kernel over \(r\in[R_j]\).
Unfolding the theta sum leaves the constant term along maps that
raise the grading, giving the next function \(\mathcal W_{j+1}\).
Thus
\[
 \begin{aligned}
 \int_{[R_j]}W_{j,g_j}(r)(L_{\xi,j}^{\mathrm t}z)(rg_j)\,dr
   &=d_{\xi,j}\,\Xi_{j+1}(z),\\
 \Xi_{j+1}(z)
   &=\int_{R_{j+1}(F)\backslash G_{j+1}(\mathbb A)}
                   \mathcal W_{j+1}(h)z(h)\,dh .
 \end{aligned}
\]
\leansrc{New/UpperBound/Seesaw/UnfoldingStages}{One step of the Bessel-period unfolding}.

At the first step, decompose the support at principal level \(J_M\)
into the classes \(\Omega_0=J_M/J_{\xi,0}\).
Unfolding the Poincar\'e series gives the average over these classes,
so the remaining factor is
\[
 \frac{d_{\xi,0}}{[J_M:J_{\xi,0}]}=[J^\circ_0:J_M],
\]
independent of the twist.
For \(1\le j<m\), a finite set \(\Omega_j\) of classes modulo
\(J_{\xi,j}\) meets the \(R_j(\mathbb A_f)\)-orbits of the finite
components on which \(\mathcal W_j\) is nonzero.
The multiplicity \(m_j(c)\) of \(c\in\Omega_j\) counts the classes
of elements of \(R_j(\mathbb A_f)\), modulo its intersection with
the standard lattice stabilizer, that move \(c\) back into the set.
Weighting by \(1/m_j(c)\) gives orbital integral \(1\).
Since \(J^\circ_j\) has volume \(1\), each \(J_{\xi,j}\)-coset has
volume \(d_{\xi,j}^{-1}\), which cancels the \(d_{\xi,j}\) in the
period. The remaining infinite integrals use the measures of this
weighted decomposition, and the finite sum runs over
\(\prod_{j<m}\Omega_j\).
The first multiplicity is \(1\), and each other multiplicity
divides a power of \(\mathrm NQ(\xi)\), the norm of the twist modulus.
\leansrc{New/UpperBound/Seesaw/IntegralAsCosetSum}{Orbital weights and the volume of the level}.
\leansrc{New/UpperBound/Seesaw/Cut}{Finite classes and their multiplicities}.

\emph{Decomposing the kernels in pairs.}
Fix the outer variables and take the product of two consecutive
constant terms for \(Y_j,Y_{j+1},Y_{j+2}\), with \(j\) even.
Reverse the first kernel by taking the adjoint of its map, and join
it with the second to obtain a datum on
\(\operatorname{Hom}(Y_{j+1},Y_j^{-}\oplus Y_{j+2})\), where \(Y_j^{-}\)
has the negative Hermitian form. If \(T,T'\) are the two maps along
the graded strings, the translation of the joined datum is
\((T^\ast,T')\).
The two means over maps that raise the grading give the projection
of the joined datum to its degree-zero part.
On the unipotent subgroup of \(R_j\), the characters of
\(\mathcal W_j\) and the constant term cancel.
Averaging over the unipotent subgroup of \(R_{j+1}\) gives the mean
over the remaining isotropic coordinates in the hyperbolic planes.
The identity \(-TT^\ast+(T')^\ast T'=0\) makes the Heisenberg phase
for joining the translations equal to \(1\).
The remaining kernel, up to a measure constant, is a theta kernel
on \(\operatorname{Hom}(X_{j+1},X_j^{-}\oplus X_{j+2})\).
Integrate first over \([U(X_{j+1})]\) and then over \([U(X_j)]\)
against the preceding matrix coefficient.
For products of Fock polynomials the kernel factors into the product
of the two definite theta kernels, so the two integrals give the
matrix coefficient of their composite lift.
Density of Fock polynomials and continuity of the integrals give the
identity for the infinite component used here.
\leansrc{New/UpperBound/Seesaw/UnfoldingStages}{Product of constant terms and the definite double integral}.

The infinite and finite factors in this calculation transform separately.
The infinite weight of the twist is fixed. The infinite factor therefore
changes only the linear functional used to take the matrix coefficient,
independently of the twist, \(\phi\), and the finite outer variables.
The finite factor gives the evaluated data \(\mathcal E_i(\xi;\boldsymbol c)\)
for each pair at the tuple of classes \(\boldsymbol c\).
Induction by pairs gives the final function \(\mathcal W_m\) as a matrix
coefficient of the lift along \(X_0\to\cdots\to X_m\) with datum
\[
  \mathcal E(\xi;\boldsymbol c)
  =\bigotimes_{i=0}^{m/2-1}\mathcal E_i(\xi;\boldsymbol c).
\]

At the top, \(Y_m\) is definite and \(R_m=U(Y_m)\).
Its point \(h\) embeds into \(X_m\) as \(\operatorname{diag}(h,1)\).
In integrating the matrix coefficient against \(\overline{x_m}\),
orthogonality at infinity gives the finite class mean after applying
the integral branching map. The inner product with the restriction to
\(Y_m\) is that same mean divided by the dimension of the Weyl module
on \(Y_m\) and by \(\prod_\tau\tau(\overline{\pi_N})^{u_\tau}\), where
\(u_\tau\) is the branching exponent.
Thus the top period is the pairing of the restricted upper lift with
\(x_m\) times a constant depending only on the outer infinite variables.
Integrating this constant and multiplying by \([J^\circ_0:J_M]\)
gives \(c_0\). It follows that \(\langle x_m,L_\xi\mathcal P\phi\rangle\)
is \(c_0\) times the pairing with the restricted upper lift for the datum
\[
 \mathcal A_\xi
   =\sum_{\boldsymbol c\in\prod_{j<m}\Omega_j}
         \left(\prod_{j<m}m_j(c_j)^{-1}\right)\mathcal E(\xi;\boldsymbol c).
\]
Since the equality holds for every test form \(x_m\), it identifies
the lifts themselves.
\leansrc{New/UpperBound/Seesaw/SeesawIdentity}{Top branching map and equality of the two chains}.

\emph{Clearing the denominators of the finite data.}
Set \(\mathsf B=M^{f n_m^2}\), with \(n_j=\dim_E X_j\).
Finite Fourier expansion expresses each evaluated datum as a sum of
glued translated data. The denominator for one pair divides
\(\mathsf B^4\), and the factor relating the vacuum to the integral
vacuum is independent of the twist.
Thus \(\mathcal A_\xi=\sum_t r_t C\,\mathcal V_t\), with \(\mathcal V_t\)
the glued translates of the integral vacuum, \(C\) independent of
the twist, and \(\mathsf B^{2m}r_t\) integral at places prime to
\(\mathrm NQ(\xi)\).
For the trivial twist the sum of the absolute values of the coefficients
of each evaluated tensor is at most \(1\), and each sum over the
classes contributes at most \(\mathsf B\). Thus
\(\sum_t|\sigma(r_t)|\le\mathsf B^m\).
The level actions permute the tuples of classes and preserve their
multiplicities. This gives joint invariance of \(\mathcal A_\xi\)
under the intermediate levels of the upper chain.
\leansrc{New/UpperBound/Seesaw/GoodDatum}{Finite Fourier expansion of the glued data}.

If \(c_0C/\mathsf B^{2m}\ne0\), take it for \(c_{\mathrm{ss}}\) and
put \(\Psi_\xi(t)=\mathsf B^{2m}r_t\).
If \(c_0C=0\), the composite map vanishes and one may take
\(c_{\mathrm{ss}}=1\) and \(\Psi_\xi=0\).
At places prime to \(M\mathrm NQ(\xi)\), the translations can be
replaced by zero with the same effect on the vacuum.
Writing \(B_{\xi,\Psi_\xi}\) for the corresponding upper composite
lift, we obtain
\begin{equation}\label{eq:upper-seesaw}
 L_\xi(c_{\mathrm{ss}}^{-1}\mathcal P\phi)=B_{\xi,\Psi_\xi}\phi
 \qquad(\phi\in S).
\end{equation}
The translations are supported above \(M\mathrm NQ(\xi)\), and
multiplication by \(\overline{\pi_M}\) makes them integral at places
prime to \(\mathrm NQ(\xi)\).
The coefficients are integral at the same places. For the trivial twist,
\(\sum_t|\sigma(\Psi_1(t))|\le\mathsf B^{3m}=M^{3fmn_m^2}\).
The glued tensor is jointly invariant under the intermediate levels.\leanref{exists_isAdmissibleCombination_seesawDatum}

For the estimates on the composite lift, expand each gluing as a
finite sum modulo \(N\). Each term adds the last column \(-y/N\) to
one translation and the last row \(-\overline y/N\) to the next.
Collecting the terms gives a sum of unglued translated data with
the same integrality condition on the coefficients.
The translations may have denominators at \(N\). For the trivial
twist the sum of the absolute values of the coefficients is at most
\(\exp(x^{60})\).
\leansrc{New/UpperBound/Divisibility/ChainDatumExpansion}{Expansion of the glued tensor}.

After projection, the lift is a sum of right translates of \(L_\xi z\).
For a principal-level form \(z\) and a finite adelic point \(u\) of \(U(Y_m)\),
\begin{equation}\label{eq:upper-project}
 L_\xi\bigl(|\overline G|e_0\rho(g)z\bigr)(u)
 =
 \sum_{v\in\overline G}|\overline G|e_0(v)\,
            (L_\xi z)\bigl(u\,\widetilde{vg}_{\,\xi}\bigr).
\end{equation}
Here \(\widetilde{vg}_{\,\xi}\) is the lift of \(vg\) to the first block at
the corresponding level. Equation \eqref{eq:upper-project} follows from
the equivariance of \(L_\xi\) for the group action.
The expression of the central projector in irreducible characters of
\(\overline G\) shows that \(|\overline G|e_0(v)\) is an algebraic integer
and \(|\sigma(e_0(v))|\le 1\).
Right translation preserves the lattices used to test integrality.
The ultrametric triangle inequality therefore shows that whenever
\(cL_\xi z\) is integral, so is \(cL_\xi(|\overline G|e_0\rho(g)z)\).
If the squared norm of each value is at most \(\exp(B)\), the projected
form has pointwise norm at most \(|\overline G|^2\exp(B/2)\).
For the trivial twist the number of classes is at most \(\exp(x^{30})\).
Squaring the pointwise estimate and summing over classes gives an upper
bound of \(\exp(B+x^{99})\).
\leansrc{New/UpperBound/Size/ArchimedeanSize}{Integrality and norm estimates for the projection}.

\subsection{Theta derivatives and contributions from the finite places}

Put \(K=\mathbb Q(\sqrt{-d})\) and normalize the valuation at \(p\)
by \(v_p(p)=1\). If \(p\) is nonsplit in \(K\), set
\[
  \pi_p=\begin{cases}p & (p\ne d),\\ \sqrt{-d}&(p=d),\end{cases}
  \qquad
  \rho_p=\begin{cases}(p^2-1)^{-1}&(p\ne d),\\
                        (2p-2)^{-1}&(p=d).\end{cases}
\]
Let \(I_p^0\subseteq\mathbb Z_{\ge0}\) be the positive odd integers if \(p\)
is inert and all nonnegative integers if \(p=d\).
Set \(I_p=I_p^0\setminus\{0\}\). At a split prime both sets are empty.
Write \(\tau_e=(-p^e)\tau\) for the action on the labels of the infinite places.
As in the formalization, the multiplier is the identity when \(-p^e\)
does not represent a unit modulo the cyclotomic modulus.

\paragraph{Normal form for the theta series}
Let \(p\) be nonsplit. Fix a \(p\)-adic embedding and write its images in
\(\mathbb C_p\). Write \(\mathcal O\) for the valuation ring and put
\(q=p^2\) in the inert case and \(q=p\) in the ramified case, so that
\(\rho_p=v_p(\pi_p)/(q-1)\). Let \(\theta\) denote the theta series of
the CM elliptic curve, and let \(y=z+O(z^2)\) be the normalized local
parameter. For the chosen point \(t\) of order two, the translated
coordinate \(x_{\rm CM}(z)=\wp(z+t)-\wp(t)\) belongs to
\(\mathcal O[[y]]\) and satisfies \(x_{\rm CM}=z^2+O(z^3)\).

\emph{Rescaling the formal parameter.}
Complex multiplication by an integer of \(K\) preserves
\(\mathcal O[[y]]\). Passing to the \(p\)-adic closure of this action gives
\(y(\zeta z)\in\mathcal O[[y]]\) for \(\zeta^{q-1}=1\). Set
\[
 T(z)=\frac{1}{q-1}\sum_{\zeta^{q-1}=1}\zeta^{-1}y(\zeta z).
\]
This is the part of \(y\) in degrees \(1\bmod(q-1)\).
Since \(q-1\) is a unit and \(T=y+O(y^2)\), the coefficients of the inverse
substitution can be solved for recursively in \(\mathcal O\).
Thus \(\mathcal O[[T]]=\mathcal O[[y]]\).
The relation \(T(\zeta z)=\zeta T(z)\) gives
\[
 T(\pi_p z)=\pi_p T(z)+\sum_{i\ge1}b_i T(z)^{1+i(q-1)}
 \qquad (b_i\in\mathcal O).
\]
The order of vanishing of this series in the residue field equals the
inseparable degree of \([\pi_p]\), which is at most its degree \(q\).
The linear term reduces to zero, and there are no terms of degrees
\(2,\ldots,q-1\). Hence \(b_1\) is a unit.
\leansrc{Lib/CM/FormalGroup/InseparableDegree}{Reduction and inseparable degree of the CM action}.

Choose \(\varpi^{q-1}=-\pi_p/b_1\) and put \(Y(z)=\varpi^{-1}T(\varpi z)\).
Then \(v_p(\varpi)=\rho_p\), and substitution in the preceding identity gives
\[
 \frac{Y(\pi_p z)}{\pi_p}=Y(z)-Y(z)^q+\pi_p H_1(Y(z)),
 \qquad H_1\in z^2\mathcal O[[z]].
\]
Indeed, the coefficient of the term indexed by \(i\ge2\) is
\(b_i\varpi^{i(q-1)}/\pi_p\in\pi_p\mathcal O\).
We have \(Y(0)=0\) and \([z]Y=1\). In degree \(n\ge2\),
\((\pi_p^{n-1}-1)[z^n]Y\) is an integral polynomial in the earlier
coefficients. Since \(\pi_p^{n-1}-1\) is a unit, induction yields
\(Y\in\mathcal O[[z]]\). Reduction in the ring \(\mathcal O/(\pi_p)\) of
characteristic \(p\) gives \(Y-Y^q=z\). There is a unique solution with
zero constant term, determined coefficient by coefficient, and it is
\begin{equation}\label{eq:upper-formalparam}
 Y(z)\equiv\sum_{\nu\ge0}z^{q^\nu}\pmod{\pi_p}.
\end{equation}
For any \(H\in\mathcal O[[y]]=\mathcal O[[T]]\), substituting
\(T(\varpi z)=\varpi Y(z)\) shows that \(H(\varpi z)\in\mathcal O[[z]]\).
Even before rescaling, the coefficients of \(H\) in degrees below \(q\)
are integral, since \(T=z+O(z^q)\).
\leansrc{Lib/CM/FormalGroup/FormalParameter}{Construction and congruence of the formal parameter}.

\emph{The exponent in one variable.}
Put \(g=-\log\theta\). Pullback by \([\pi_p]\) makes \(\theta\) an even
section of degree \(q\) and preserves integrality at torsion points.
The expansion of even sections gives
\(\theta(\pi_p z)=\theta(z)^qR(x_{\rm CM}(z))\) with \(2\deg R<q\),
where the inequality is strict because \(q\) is odd.
\leansrc{Lib/CM/Integrality/IntegralSections}{Polynomial expansion of even sections}.
This identity gives
\begin{equation}\label{eq:upper-frobenius}
 q g(z)-g(\pi_p z)=\pi_p H_2,\qquad H_2\in\mathcal O[[y]].
\end{equation}
To see this, write \(g_n=[z^n]g\). We have \(g_0=g_1=0\) and
\(\pi_p g_n\in\mathcal O\) for \(n<q\).
For \(n=2\), this follows from the integrality of \(\sqrt{-d}\,g_2\)
and \(v_p(\sqrt{-d})\le v_p(\pi_p)\).
\leansrc{Lib/CM/Integrality/LogTheta}{Integrality of the quadratic coefficient of the theta logarithm}.
For \(n>2\), the identity \(g''=g''(0)+x_{\rm CM}\) gives
\(n(n-1)g_n=[z^{n-2}]x_{\rm CM}\). The right-hand side is integral
if \(n<q\), and then \(v_p(n(n-1))\le1\) in the inert case and
\(v_p(n(n-1))=0\) in the ramified case.
Since \(q=\pi_p^2\) or \(q=-\pi_p^2\),
\((q-\pi_p^n)g_n\in\pi_p\mathcal O\) in degrees below \(q\).

It follows that all coefficients of \(\exp(qg-g(\pi_p z))-1\) in degrees
below \(q\) are divisible by \(\pi_p\). Indeed,
\(v_p(j!)\le(j-1)/(p-1)\le (j-1)v_p(\pi_p)\) for \(j\ge1\), and the
exponent has constant term zero. The exponential is \(R(x_{\rm CM})\).
The first term of \(x_{\rm CM}^i\) is \(z^{2i}\), so comparison of the
coefficients in degree \(2i\), starting at \(i=0\), yields
\(R-1\in\pi_p\mathcal O[X]\). Hence \(R(x_{\rm CM})=1+\pi_p H\) with
\(H\in\mathcal O[[y]]\) and \(H(0)=0\).
The logarithm belongs to \(\pi_p\mathcal O[[y]]\), since its \(j\)-th term
has scalar \(\pi_p^j/j\in\pi_p\mathcal O\).
\leansrc{New/Derivatives/FormalGroup/FrobeniusCongruence}{Frobenius congruence for the theta logarithm}.

Let \(g^\flat\) be the part in degrees \(p+1\bmod(q-1)\).
Write \(i_0\) for the least nonnegative representative of this class.
It equals \(p+1\) in the inert case and \(2\) in the ramified case.
The integrality of the theta transformation under the CM action gives
\[
 \exp\bigl(\zeta^{-(p+1)}g(\zeta z)-g(z)\bigr)\in\mathcal O[[y]]
 \qquad(\zeta^{q-1}=1).
\]
This follows by passing to the limit of the transformation at CM integers
and their norms, using the pair \((\zeta^{-1},\zeta^{-(p+1)})\), and
replacing the variable by \(\zeta z\).
\leansrc{Lib/CM/FormalGroup/FormalParameter}{Theta transformation at roots of unity}.
Multiplying these identities over all \(\zeta\) gives
\(\exp(g^\flat-g)^{q-1}\in\mathcal O[[y]]\) by the averaging formula
for this degree class. The \((q-1)\)-st root with constant term \(1\)
is obtained recursively by dividing the coefficient at each new degree
by \(q-1\). Hence \(\exp(g^\flat-g)\) itself belongs to \(\mathcal O[[y]]\).
With \(c_\varpi=\varpi^{p+1}/p\), this gives
\begin{equation}\label{eq:upper-one-variable}
 \begin{aligned}
 \theta(\varpi z)&=W_1(z)\exp(c_\varpi\Xi(z)),\\
 W_1(z)&=\exp\bigl((g^\flat-g)(\varpi z)\bigr)\in\mathcal O[[z]],\\
 \Xi(z)&=-g^\flat(\varpi z)/c_\varpi .
 \end{aligned}
\end{equation}

Take the part of \eqref{eq:upper-frobenius} in degrees
\(i_0\bmod(q-1)\). Since \(T(\zeta z)=\zeta T(z)\), one may take the
corresponding degree class directly in the expression for \(H_2\) in
\(\mathcal O[[T]]\). This part is
\(\lambda' T^{i_0}+T^{i_0+q-1}H_3\) with
\(\lambda'\in\mathcal O\) and \(H_3\in\mathcal O[[T]]\).
Replace \(z\) by \(\varpi z\) and divide by \(-q c_\varpi\).
The scalar \(-\pi_p\varpi^{i_0}/(q c_\varpi)\) is a unit, and the
second term retains a factor \(\varpi^{q-1}\in\pi_p\mathcal O\).
For some \(\lambda\in\mathcal O\) we thus obtain
\[
 (1-\pi_p^n/q)[z^n]\Xi\equiv\lambda[z^n]Y^{i_0}\pmod{\pi_p\mathcal O}.
\]
The coefficients for \(n<i_0\) vanish. If \(n>2\),
\(1-\pi_p^n/q\) is a unit congruent to \(1\) modulo \(\pi_p\).
For \(n=2\) in the ramified case, this factor equals \(2\).
Thus \(\Xi\) is integral, with reduction
\[
 \Xi\equiv
 \begin{cases}
    \lambda Y^{p+1}& (p\text{ inert}),\\
    \lambda(Y^2-z^2/2)& (p=d)
 \end{cases}
 \pmod{\pi_p}.
\]
In the inert case, expanding \(Y^pY\) using \eqref{eq:upper-formalparam}
gives exponents \(p^{2i+1}+p^{2j}\). The two exponents of \(p\) have
positive odd difference. Each pair consisting of the smaller exponent
\(\nu\) and the difference \(e\) occurs once. In the ramified case,
the diagonal terms of \(Y^2\) have coefficient \(1\), and the terms
with distinct exponents have coefficient \(2\). The correction changes
the coefficient of \(z^2\) to \(1/2\).
This gives
\[
  \Xi(z)\equiv
  \sum_{e\in I_p^0}\sum_{\nu\ge0}\lambda_{e,\nu}(z^{p^e}z)^{p^\nu}
  \pmod{\pi_p}.
\]
In the inert case all coefficients \(\lambda_{e,\nu}\) equal \(\lambda\).
In the ramified case \(\lambda_{0,0}=\lambda/2\),
\(\lambda_{0,\nu}=\lambda\) for \(\nu>0\), and \(\lambda_{e,\nu}=2\lambda\)
for \(e>0\).
\leansrc{New/Derivatives/FormalGroup/Exponent}{Construction and coefficient congruence of the one-variable exponent}.

\emph{The multivariable theta product.}
Suppose that the vectors \(u_i\) are integral at all primes above \(p\)
and that the multiplicities \(\mu_i\in\mathbb Z_{\ge0}\) satisfy
\(\sum_i\mu_i u_i u_i^\ast=1\).
Write \(L_i(z)=u_i^\ast z=\sum_{\tau,c}\overline{\tau(u_{i,c})}z_{\tau,c}\)
for the full linear form. Let \(F\) be the product of
\(\prod_i\theta(L_i(z))^{\mu_i}\) and an integral power series in
the \(y(L_i(z))\). Substitute each \(L_i\) in \eqref{eq:upper-one-variable}.
To simplify the exponent, use
\(\overline{\tau(u_{i,c})}^{\,p^e}\equiv\tau_e(u_{i,c})\pmod{\pi_p}\).
For algebraic integers, this congruence follows by comparing Frobenius
and the partner action on the generating root of unity and on
\(\sqrt{-d}\). At an inert prime, \(e\) is odd, so conjugation and the
\(p^e\)-th power act with the same sign on \(\sqrt{-d}\).
At a ramified prime \(\sqrt{-d}\equiv0\pmod{\pi_p}\).
For a vector entry integral above \(p\), one can clear denominators by a
rational integer prime to \(p\), whose \(p^e\)-th power is congruent to
itself. Thus
\[
 \begin{aligned}
 \sum_i\mu_i L_i(z)^{p^e}L_i(z)
 &\equiv \sum_{\tau,c}z_{\tau,c}^{p^e}
              \sum_i\mu_i\tau_e(u_{i,c}) L_i(z)\\
 &=\sum_{\tau,c} z_{\tau,c}^{p^e}z_{\tau_e,c}\pmod{\pi_p}.
 \end{aligned}
\]
The last equality is the \((\tau_e,c)\) coordinate of the resolution
of the identity \(\sum_i\mu_i u_i u_i^\ast=1\).
Moreover, taking \(p^\nu\)-th powers commutes with summation after
reduction, and \(\mu_i^{p^\nu}\equiv\mu_i\).
Consequently \(\sum_i\mu_i\Xi(L_i(z))\) differs from the sum in the
following exponent by an element of \(\pi_p\mathcal O[[z]]\) with zero
constant term.
\begin{equation}\label{eq:upper-normal}
 F(\varpi z)=W(z)\exp\!\left(
       \frac{\varpi^{p+1}}{p}
       \sum_{e\in I_p^0}\sum_{\nu\ge0}\lambda_{e,\nu}\Gamma_e(z)^{p^\nu}
                              \right),
 \qquad
 \Gamma_e(z)=\sum_{\tau,c}z_{\tau,c}^{p^e}z_{\tau_e,c}.
\end{equation}
Here \(W\in\mathcal O[[z]]\). Indeed, the exponential of the error
multiplied by \(c_\varpi\) is integral because
\(v_p(c_\varpi\pi_p)=1/(p-1)\) and \(v_p(j!)\le j/(p-1)\).
Multiplying this factor by \(\prod_i W_1(L_i(z))^{\mu_i}\)
and the integral series in the \(y(L_i(z))\) rescaled at \(\varpi z\)
gives \(W\). Each coefficient of the exponent receives only finitely
many contributions.\leanref{exists_hasNormalForm}

\paragraph{Differentiation along products of minors}
For a polynomial \(Q\), write \(J_Q(F)=(Q(\partial)F)(0)\).
Suppose \(Q\) is homogeneous of degree \(B\) with integral coefficients and
that all coefficients of \(G\) of degree \(B\) are integral.
If there are \(w\) variables, then
\begin{equation}\label{eq:upper-legendre}
  v_p(J_Q(G))\ge
  \frac{B}{p-1}-w\bigl(1+\lfloor\log_p\max(B,1)\rfloor\bigr).
\end{equation}
Indeed, the derivative for a monomial of degree \(|D|=B\) multiplies
its coefficient by \(D!=\prod_c D_c!\). Legendre's formula gives
\(v_p(D_c!)=(D_c-s_p^{\mathrm{dig}}(D_c))/(p-1)\), where
\(s_p^{\mathrm{dig}}\) is the sum of the base-\(p\) digits.
Each digit is at most \(p-1\) and there are at most
\(1+\lfloor\log_p\max(B,1)\rfloor\) digits. Apply the ultrametric
triangle inequality when summing the terms.

\emph{From a Schur matrix entry to a product of minors.}
Write \(n=n_j,\ n'=n_{j+1}\) for the source and target ranks in one step.
The diagram at each place \(\tau\) has no column of height strictly between
\(n\) and \(n'\). Fix a standard filling \(I'\) at the target. Let
\(Z_\tau=(z_{\tau,b,a})\) be an \(n'\)-by-\(n\) matrix of variables,
and \(Y_\tau=(y_{\tau,a,k})\) a matrix of integral directions.
Let \(Q_Y(Z)\) be the product of minors of \(Z_\tau Y_\tau\) obtained
by taking, for each column of height \(h<n'\) at \(\tau\), the \(h\) rows
specified by \(I'\) and the first \(h\) columns.
Write \(P_{I'I}(Z)\) for the \((I',I)\) entry of the Schur matrix.
Its defining expansion in the standard fillings gives
\[
 Q_Y(Z)=\sum_I B_I(Y)P_{I'I}(Z).
\]
Here \(B_I(Y)\) is the product of minors for the source filling \(I\)
times a common determinant factor from columns of height equal to
the source rank when \(n<n'\). The determinant factor is a nonzero
polynomial over the residue field, and the standard-filling basis
is linearly independent over that field. Hence the reductions of the
\(B_I\) are linearly independent.
\leansrc{New/Derivatives/Divisibility/Count}{Expansion of the Schur matrix in the standard-filling basis}.

Put \(\alpha_I=J_{P_{I'I}}(F)\). If all \(\alpha_I\) vanish, the bound is
immediate. Otherwise choose \(|\alpha_{I_0}|_p\) maximal. The reduction
of \(\sum_I(\alpha_I/\alpha_{I_0})B_I\) is a nonzero polynomial.
The residue field is infinite, so choose a point at which the reduction
is nonzero and lift its coordinates to \(\mathcal O\).
At the resulting \(Y\), the norm of
\(J_{Q_Y}(F)=\sum_I\alpha_I B_I(Y)\) is \(|\alpha_{I_0}|_p\).
A lower bound on \(v_p(J_{Q_Y}(F))\) for all integral \(Y\) thus
gives the same lower bound for every Schur derivative entry.\leanref{exists_norm_schurJet_le_norm_jet_minorProd}

\emph{Congruences on the directions.}
Let \(N_{\tau h}\) be the number of columns of height \(h<n'\).
Write \(t_\tau\) for the number of columns of height \(1\),
\(r_\tau\) for the number of the remaining columns, and \(B_\tau\) for
their number of boxes. The total number of boxes is
\(B=\sum_\tau(t_\tau+B_\tau)\). Let \(A\) be the set of marked places.
In counting congruences, put \(t_\tau(A)=0\) at \(\tau\in A\) and
\(t_\tau(A)=t_\tau\) elsewhere.
At each \(\tau,h\), an \(h\)-by-\(h\) matrix \(g_{\tau h}\) changes the
directions to
\[
 \beta_{\tau h k,a}=\sum_{\ell<h}(g_{\tau h})_{k\ell}\,y_{\tau,a,\ell}.
\]
For \(e\in I_p^0\) with \(e\le\lfloor\log_p\max(B,1)\rfloor\) and
\(\tau'=\tau_e\), the required congruence is
\begin{equation}\label{eq:upper-direction-pairing}
 \sum_a\beta_{\tau h k,a}^{\,p^e}\beta_{\tau' h' k',a}\equiv 0
   \pmod{\pi_p}.
\end{equation}
A pair of heights is exempt when \(h=1\) and either \(\tau\in A\)
or both \(h'=1\) and \(\tau'\in A\).

The conditions on the rows of \(g_{\tau h}\) are linear.
Choose \(\alpha_e^{p^e}=\pi_p\). Use modulus \(\alpha_e\) on the left
and modulus \(\pi_p\) on the right. If every entry of a row on the
left is divisible by \(\alpha_e\), each component of the new direction
has \(p^e\)-th power divisible by \(\pi_p\).
On the right, require each entry of the row to be divisible by \(\pi_p\).
If the height on the left is \(1\), its new direction is a scalar
multiple of the original direction \(y_{\tau,\cdot,0}\).
For each row on the right it suffices to impose one congruence,
\(\sum_a y_{\tau,a,0}^{p^e}\beta_{\tau' h' k',a}\equiv0\pmod{\pi_p}\).
If the height on the right is \(1\), choose a \(p^e\)-th root
\(\eta_a\in\mathcal O\) of each \(y_{\tau',a,0}\) and impose
\(\sum_a\eta_a\beta_{\tau h k,a}\equiv0\pmod{\alpha_e}\).
Raising to the \(p^e\)-th power gives
\(\sum_a y_{\tau',a,0}\beta_{\tau h k,a}^{p^e}\equiv0\pmod{\pi_p}\),
since \(\pi_p\mid p\) makes the \(p^e\)-th power additive after reduction.
Thus a pairing condition uses only the original direction on the
opposite side.

Allocate the conditions as follows.

\begin{enumerate}
\item If the height on the left is \(1\), impose entrywise conditions
at height \(1\) on the left, or impose pairing conditions on the right
against the original direction on the left. The respective costs are
\(t_\tau(A)/p^e\) and \(r_{\tau'}+t_{\tau'}(A)\). Choose the smaller cost.
Omit marked columns of height \(1\) on the chosen side.
\item If the height on the left exceeds \(1\) and the height on the
right is \(1\), impose pairing conditions on the left against the
original direction on the right, at cost \(r_\tau/p^e\).
\item If both heights exceed \(1\), impose entrywise conditions on the
relevant rows on one side, taking the smaller of the costs
\(B_\tau/p^e\) and \(B_{\tau'}\).
\end{enumerate}

The unit of cost is \(v_p(\pi_p)\). At each height a pairing condition
is one linear form, and an entrywise condition gives as many linear
forms as the height. The modulus on the left contributes
\(v_p(\alpha_e)=v_p(\pi_p)/p^e\), accounting for the factor \(p^{-e}\).
The factor given by the height counts the linear forms for the
entrywise conditions. Requiring divisibility of one form on every
row costs only one copy of its modulus in the determinant.

Indeed, let \(a_k\) be the values of a linear form on the rows.
If some value is nonzero, choose \(a_{k_0}\) with minimal valuation
and subtract \(a_k/a_{k_0}\) times row \(k_0\) from each other row.
The quotient is integral, and the values on those rows become zero.
Multiply the remaining row by the modulus. This multiplies the
determinant by that modulus. If all values are zero, the identity
operation suffices. Starting from the identity matrix, impose the
conditions one form at a time. The new rows are integral linear
combinations of the preceding ones, so every earlier divisibility
condition is preserved. Thus \(\det g_{\tau h}\) divides the product
of the moduli imposed at its place and height. The same matrix occurs
once for each column of that height, so
\[
 \begin{aligned}
  \sum_{\tau}\sum_{h<n'}N_{\tau h}v_p(\det g_{\tau h})
   &\le v_p(\pi_p)\mathcal D_p(A),\\
  \mathcal D_p(A)
   &=\sum_{\substack{e\in I_p^0\\e\le\lfloor\log_p\max(B,1)\rfloor}}
       \sum_{\tau\in\Phi}
       \biggl\{\min\!\left(\frac{t_\tau(A)}{p^e},\,r_{\tau_e}+t_{\tau_e}(A)\right)\\
   &\hphantom{=\sum_e\sum_\tau\biggl\{}{}
          +\frac{r_\tau}{p^e}
          +\min\!\left(\frac{B_\tau}{p^e},\,B_{\tau_e}\right)\biggr\}.
 \end{aligned}
\]
This establishes \eqref{eq:upper-direction-pairing} for all the pairs
that are subject to it, simultaneously.\leanref{exists_newDirection_dvd}

\emph{The valuation bound by the chain rule.}
Introduce variables \(s_{\tau h k b}\) and put
\[
 z_{\tau,b,a}(s)=\sum_{h<n'}\sum_{k<h}s_{\tau h k b}\,p^{w_{\tau h b}}\beta_{\tau h k,a}.
\]
When no places are marked, put \(w_{\tau h b}=0\).
When places are marked, choose nonnegative integers \(w_{\tau h b}\)
so that at least one of \(w_{\tau 1 b}\) and \(w_{\tau' h' b}\) is
positive for each exempt pair and each row \(b\).
Then \(\Gamma_e(z(s))\) is divisible by \(\pi_p\). Indeed, after
distributing the \(p^e\)-th power over the sum in the quotient ring,
each term has coefficient equal to the pairing in
\eqref{eq:upper-direction-pairing} times
\(p^{p^e w_{\tau h b}+w_{\tau' h' b}}\).
The exempt pairs also vanish by this factor.

Let \(Q_*(s)\) be the product of the determinants
\(\det(s_{\tau h k b})\) on the rows \(b\) chosen by \(I'\) and
directions \(k<h\) for each column. Let \(W_{\rm ch}\) be the sum of
\(w_{\tau h b}\) over the boxes, using the row index \(b\) of each box.
The chain rule gives
\begin{equation}\label{eq:upper-minor-chain-rule}
 J_{Q_*}\bigl(F(\varpi z(s))\bigr)
 = \varpi^B p^{W_{\rm ch}}\!
    \left(\prod_{\tau,h<n'}(\det g_{\tau h})^{N_{\tau h}}\right)J_{Q_Y}(F).
\end{equation}
In one column, the derivative matrix along the new directions is the
derivative matrix for \(Y_\tau\) multiplied on the right by
\(g_{\tau h}^{\,t}\) and on the left by the diagonal matrix with row
factors \(\varpi p^{w_{\tau h b}}\). Taking the product of determinants
gives \eqref{eq:upper-minor-chain-rule}.

The polynomial \(Q_*\) is integral of degree \(B\) in at most
\(f(n')^3\) variables. In \eqref{eq:upper-normal} after substitution,
the terms with \(e>\lfloor\log_p\max(B,1)\rfloor\) have degree
greater than \(B\). All remaining couplings are divisible by \(\pi_p\),
so the exponential calculation above shows that \(F(\varpi z(s))\)
has integral coefficients through degree \(B\).
Equations \eqref{eq:upper-legendre} and \eqref{eq:upper-minor-chain-rule} give
\begin{equation}\label{eq:upper-jet}
 \begin{aligned}
 v_p(J_{Q_Y}(F))+W_{\rm ch}
   &\ge \gamma_p B-v_p(\pi_p)\mathcal D_p(A)\\
   &\quad{}-f n_{j+1}^3\bigl(1+\lfloor\log_p\max(B,1)\rfloor\bigr),\\
 \gamma_p&=\frac1{p-1}-\rho_p .
 \end{aligned}.
\end{equation}
Here \(B\) includes the boxes of marked columns.\leanref{norm_jet_minorProd_le_of_charges}
For the positive exponents one has
\(\gamma_p=v_p(\pi_p)\sum_{e\in I_p}p^{-e}\).
Subtracting the deductions for \(e>0\) in \(\mathcal D_p(A)\) from
\(\gamma_p B\) gives the following contributions in units of \(v_p(\pi_p)\).
\[
 \left(\frac{t_\tau(A)}{p^e}-t_{\tau_e}(A)-r_{\tau_e}\right)_+
 +\left(\frac{B_\tau}{p^e}-B_{\tau_e}\right)_+
 -\frac{r_\tau}{p^e}.
\]
This uses \(x-\min(x,y)=(x-y)_+\) for the two minima.
Extending the deductions to all of \(I_p\) still gives a lower bound,
since the additional terms are nonnegative.
Add \(\gamma_p\) for each marked box. At \(p=d\), also subtract the
deduction for \(e=0\).

\paragraph{Gain in one step}
Define the local prime sum and the gain over the full chain by
\begin{align}
 s_p(y,\tau)
  &=v_p(\pi_p)\sum_{e\in I_p}
       \left(\frac{y}{p^e}-\kappa((-p^e)\tau)\right)_+,
       \label{eq:upper-sp}\\
 G_p
  &=m\left\{\sum_{\tau\in\Phi_a}s_p(a,\tau)
           +k\sum_{\tau\in\Phi}s_p(\kappa(\tau),\tau)\right\}.
       \label{eq:upper-gain}
\end{align}
Set \(s_p=G_p=0\) at split primes.
For \(0\le y\le a\), all the positive parts also vanish if \(p\ge x_1\).
If \(p\) is inert, \(1-\chi_K(p^e)\) is \(2\) for odd \(e\) and \(0\)
for even \(e\). If \(p=d\), this coefficient is \(1\) and
\(v_p(\pi_p)=1/2\). Grouping the terms by prime powers gives
\begin{equation}\label{eq:upper-gain-sum}
  2\sum_p G_p\log p =
    m\sum_{\tau\in\Phi_a}s_{\mathrm{fin}}(a,\tau)
       +mk\sum_{\tau\in\Phi}s_{\mathrm{fin}}(\kappa(\tau),\tau).
\end{equation}

There are \(a\) columns of height \(1\) at places in \(\Phi_a\) and none
elsewhere. The number of remaining columns is at most
\(\kappa(\tau)+x^{60}\), and the number of their boxes is between
\(k\kappa(\tau)\) and \(k\kappa(\tau)+x^{60}\).
Insert these bounds into the positive parts above.
Write \(D_{p,j}(A)\) for the valuation bound for one step, where \(A\subseteq\Phi\)
is the set of places at which the height-\(1\) columns are marked.
It is the predicted valuation of the diagram with those columns removed,
plus \(\gamma_p\) for each box removed, less the derivative error
\(f n_{j+1}^3(1+\log_p\max(B,1))\). When \(p<x_1\) is nonsplit,
\begin{equation}\label{eq:upper-one-step}
 D_{p,j}(A)\ \ge\ \frac{G_p}{m}
  - C_p(A)
  +v_p(\pi_p)|\Phi_a\cap A|\,\frac{a}{x_1}
  -\frac{fa(\log b)^2}{mp},
 \qquad
 C_p(A)=\sum_{\substack{\tau\in\Phi,\ e\in I_p\\
                 \tau,\,(-p^e)\tau\in\Phi_a\setminus A}}\frac{a}{p^e}.
\end{equation}
When a column of height \(1\) remains at the partner place, the overlap
costs at most \(a/p^e\).
The error in the other columns is controlled by \(x^{60}\) as above.
The deduction for \(p=d,e=0\) is bounded using that \(\tau,-\tau\)
cannot both lie in \(\Phi_a\) and cannot both lie in \(\Phi_b\).
In units of \(a/(mp)\), the error per place is at most \(2\) from
comparing the positive parts, \(2\) from the ramified-prime deduction,
and \(1\) from the logarithmic error in differentiation.
The condition \(5\le(\log b)^2\) absorbs these errors.
\leanref{le_divisibilityBound}

\paragraph{The local kernel and its lattice}
The local kernel of a lift \(X_j\to X_{j+1}\) uses Schwartz data on
the standard space of \(n'\)-by-\(n\) matrices.
Standard data of level \(1\) are finite \(p\)-integral combinations
of translations of the integral vacuum by points integral above \(p\).
A nonsplit prime \(p\) is prime to the cyclotomic modulus \(q_1\).
We show that the theta series of any standard datum of level \(1\)
has the normal form \eqref{eq:upper-normal}.

\emph{The theta product for the vacuum.}
Choose elements \(v_i\) of the self-dual lattice \(\mathfrak a_0\) and an
integer \(b_0\) prime to \(p\) with
\(\sum_i \tau(v_i)\overline{\sigma(v_i)}=b_0^2\delta_{\tau,\sigma}\).
The tight-frame construction makes a product of copies of \(\theta\)
along the linear forms given by placing \(v_i/b_0\) in each coordinate
in turn. The associated Weil vector \(F_0\) has theta value \(1\),
and its theta function has the prescribed signed periods along
\(b_0\mathfrak a_0\) in each coordinate.
\leansrc{New/Derivatives/Vacuum/TightFrames}{Construction of the theta product from a tight frame}.

Apply finite Fourier inversion for this sublattice to express the
integral vacuum as a sum of translates of \(F_0\).
For a representative \(u\) of a residue class, the corresponding
Fourier coefficient times the vacuum equals the sum of translates
of \(F_0\), first by \(-u\) and then by the representatives of the
period lattice. Applying the theta functional to the inversion
identity shows that the sum of these Fourier coefficients weighted
by the corresponding theta values of translates of the vacuum is
\((b_0^2)^{f n n'}\). Thus the integral vacuum is a finite sum of
translates of \(F_0\). The denominator in the coefficients is
\((b_0^2)^{f n n'}\), and the numerators are products of values of the
integral vacuum and the scalar factors of the Heisenberg action.
The coefficients are therefore integral at \(p\).
The translation points lie in the standard lattice after
multiplication by \(b_0\).
\leansrc{New/Derivatives/Vacuum/StandardDataProducts}{Finite Fourier inversion for the vacuum and its denominators}.

Apply this expansion to each translate in the standard datum.
The composition law for the Heisenberg action adds the translation
points and multiplies by an integral phase.
Each resulting translate of \(F_0\) can be represented by a
rational point \(\gamma\) inducing the same action on \(F_0\),
with \(a_0\gamma\) in the lattice in each coordinate for some
integer \(a_0\) prime to \(p\).
\leansrc{New/Derivatives/Vacuum/StandardDataProducts}{Choosing rational translation representatives}.
Each shift \(t_i\) in the theta product for \(\gamma\) satisfies
\(a_0 b_0 t_i\in\Omega_0\mathcal O_K\).
Indeed, the shift is \(-\Omega_0/b_0\) times the Hermitian pairing
between \(v_i\) and the corresponding coordinate of \(\gamma\).
The pairing of \(v_i\) and \(a_0\gamma\) lies in \(\mathcal O_K\)
by self-duality. Choose a common multiple \(k_0\) of the integers
\(a_0 b_0\) for all these translates, still prime to \(p\).

Place \(v_i/(b_0 k_0)\) in each coordinate in turn to obtain
the vectors \(u_i\) with multiplicity \(k_0^2\). Dividing the
frame identity by \(b_0^2\) gives \(\sum_i k_0^2u_i u_i^\ast=1\).
Since \(k_0 t_i\) is a period, the one-variable translated theta
series after pullback by \([k_0]\) is the product of \(\theta(z)^{k_0^2}\)
and an element of \(\mathcal O[[y]]\). Here the pulled-back
section has degree \(k_0^2\).
\leansrc{Lib/CM/Integrality/IsogenyIntegrality}{Pullback of the theta series translated by a torsion point}.
Since \(k_0\) was chosen simultaneously for all terms, the finite sum
has the common factor \(\prod_i\theta(u_i^\ast z)^{k_0^2}\).
The remaining factor lies in the ring of integral series in the
\(y(u_i^\ast z)\).

The construction first gives integrality of the adelic image
\(\alpha^{-1}u_i\). The idele \(\alpha\) that makes the standard
lattice self-dual satisfies \(q_1\alpha\) integral, and is therefore
integral at the places above \(p\) since \(p\nmid q_1\).
Thus \(u_i=\alpha(\alpha^{-1}u_i)\) is integral at all places above \(p\).
The expansion applies to the theta function with variable scaled by
\(\sqrt{2\pi/\sqrt d}/\Omega_0\). The kernel uses an additional
period factor and a factor \(\pi_M\), with \(\pi_M\overline{\pi_M}=M\).
The period factor is a unit at \(p\) since \(p\nmid q_1\).
If \(c\) is the product of the two additional factors, \(v_p(c)\ge0\).
Replacing \(z\) by \(cz\) in \eqref{eq:upper-normal} preserves the
integrality of \(W(cz)\) and multiplies the exponent coefficients by
\(c^{(p^e+1)p^\nu}\). The kernel therefore has the same normal form
at its prescribed scale.
\leansrc{New/Derivatives/Vacuum/Divisibility}{Normalization and integrality of the theta series of the local kernel}.

\emph{The bound in lattice coordinates.}
Use the standard-filling basis of each Weyl module.
The local norm \(\|\eta\|_g\) of a vector \(\eta\) at a point \(g\)
is the maximum \(p\)-adic absolute value of its coordinates after
acting by \(g^{-1}\) in the weight representation.
The weight action at \(\tau\) is read through the embedding
\(x\mapsto\iota_p(\tau(\overline x))\). The norm is unchanged by
right multiplication of the point by a lattice-preserving element,
or by the simultaneous action of a rational point on \(g\) and \(\eta\).

For a standard datum \(\psi\), write \(K_j(\psi;g,h)\) for the
matrix of Schur derivatives of \(\omega_j(g,h)\psi\) under the
twisted Weil action. Here \(g\) and \(h\) are finite adelic points
of \(X_j\) and \(X_{j+1}\), respectively.
Consider also replacing \(\psi\) by \(\omega_j(1,x d_\ast y)\psi\),
where \(x,y\) are integral above \(p\) and \(d_\ast\) is read at
each weight embedding as a diagonal matrix with entries \(u_{\tau,b}\).
Weak approximation gives rational points \(\gamma,\epsilon\) such
that \(g_0=\gamma g,\ h_0=\epsilon h x d_\ast\), and
\(d_\ast h_0 d_\ast^{-1}\) are integral above \(p\).
The points \(g_0,h_0 y\) preserve the standard lattices and their
Weil action preserves the span of integral translates of the vacuum.
The character factors of the twist are units at these points.
Thus \(\psi'=\omega_j(g_0,h_0 y)\psi\) is standard.

By covariance of the kernel and invariance of the local norms,
the comparison at the input point \(g\) and output point \(h\)
reduces to the derivative matrix of the standard datum \(\psi'\)
at the identity, with input norm at \(1\) and output norm at
\(d_\ast^{-1}\). Indeed, \(g_0\) is integral and the last two
factors of \(\epsilon h=d_\ast^{-1}(d_\ast h_0d_\ast^{-1})x^{-1}\)
are integral. The row indexed by \(I'\) is now multiplied in
absolute value by \(|u_{\tau,b}|_p\) for each box whose entry at
\(\tau\) is \(b\), and by the determinant factor
\(\prod_\tau(\prod_b|u_{\tau,b}|_p)^{e'_\tau}\), where \(e'_\tau\)
is the last entry of the target weight.
The operator norm is bounded by the largest absolute value of an
entry, so it suffices to bound these factors times the corresponding
Schur derivative.
\leansrc{New/UpperBound/Kernels/ThetaLiftGauge}{Twisted Weil action and the kernel bound in lattice coordinates}.

For unmarked standard data, take \(d_\ast=1\), so the row factors
are all \(1\). In \eqref{eq:upper-jet}, put \(W_{\rm ch}=0\),
extend the deductions to all \(e\), and relax the error term by
replacing the floor of the logarithm by the logarithm.
Each entry has valuation at least \(D_{p,j}(\varnothing)\).
Consequently
\begin{equation}\label{eq:upper-kernel-gauge}
 \|K_j(\psi;g,h)\eta\|_h
 \le p^{-D_{p,j}(\varnothing)}\|\eta\|_g
\end{equation}
for all \(g,h,\eta\).
\leansrc{New/UpperBound/Divisibility/ValuationAtNonSplitPrimes}{Valuation of the kernel for one lift}.

\emph{Adding valuations through the chain.}
Let \(\psi_{t,j}\) be the datum in step \(j\), and suppose its kernel
multiplies local norms by at most \(p^{c_j}\), uniformly for
\(t\) in the support of \(\Psi\).
Joint invariance at the intermediate levels expresses the value
of the composite lift as the finite sum, over \(t\) and class
representatives \(g_j\), of products
\[
  \Psi(t)\,K_{m-1}(\psi_{t,m-1};g_{m-1},g_m)
        \cdots K_0(\psi_{t,0};g_0,g_1)\phi(g_0).
\]
This expression has no denominator from the order of a class stabilizer.
The output lattice at \(g_{j+1}\) is the input lattice of the next
kernel, so the operator bounds multiply. If \(\phi\) has local norm
at most \(1\) and the coefficients \(\Psi(t)\) are \(p\)-integral,
the ultrametric triangle inequality bounds the sum by
\(p^{\sum_{j<m}c_j}\). Applying linear operators that commute with
the level actions to the data preserves joint invariance.
The final restriction to \(Y_m\subset X_m\) uses an integral branching
matrix and the \(p\)-adic unit \(\overline{\pi_N}\) when \(p\ne N\).
It preserves the bound.
\leansrc{New/UpperBound/Divisibility/ChainValuation}{Finite class sum and restriction bounds for the chain}.

\subsection{Colength at the finite places}

Let \(\iota\) be the isomorphism defining a valuation above \(p\), and let
\(\mathcal O_\iota\) be its valuation ring. Write \(M_\iota\) for the
module of \(\iota\)-integral forms in \(S_0\).
We say that a submodule has colength dividing \(p^j\) if for every finite
family in \(M_\iota\) a square matrix over \(\mathcal O_\iota\) sends that
family into the submodule and has determinant dividing \(p^j\).
In particular, the valuation of the determinant is at most \(j\).
This definition works directly with finite families in the module.

Let \(\mathcal U\) be the set of nonsplit primes below \(x_1\), and put
\[
 \ell_p=\frac{fa(\log b)^2}{p}
      \begin{cases} m & (p\equiv\pm1\bmod m),\\1 & (\text{otherwise}) \end{cases},
 \qquad
 \varepsilon_p=
      \begin{cases}
       2\ell_p &(p\in\mathcal U),\\
       2mfb &(p=N),\\
       0 &(\text{otherwise}).
      \end{cases}
\]
Here \(N\ge x_1\). For each \(p,\iota\), we will show that there is a
submodule \(M'_\iota\subseteq M_\iota\cap S\) of colength dividing
\(p^{j_\iota}\) in \(M_\iota\cap S\), with
\begin{equation}\label{eq:upper-local}
 j_\iota\le\varepsilon_p\dim S_0,\qquad
 v_p(B_{\xi,\Psi}\phi)\ge G_p-\varepsilon_p
  \quad(\phi\in M'_\iota,\ (p,\mathrm NQ(\xi))=1).
\end{equation}
For a form, a lower valuation bound \(v\) means that all its values have
\(p\)-adic norm at most \(p^{-v}\) in the corresponding lattice coordinates.
Here \(\Psi\) is any datum satisfying the conditions imposed in
\eqref{eq:upper-seesaw}. The choice of \(M'_\iota\) is independent of
\(\xi,\Psi\).

For \(p\notin\mathcal U\cup\{N\}\), one has \(G_p=0\).
Since \(d<x_1\), we also have \(p\ne d\), so the kernel at each step
relative to the standard lattice and the final restriction both preserve
integrality.
For \(p\in\mathcal U\), first apply \eqref{eq:upper-one-step} at every
step with no marks. The transformation laws of the kernels allow the
lower valuation bounds to be added along the chain. The coefficients
of the composite lift are integral, so its finite sum preserves the
resulting bound. This gives \(G_p-\ell_p-mC_p(\varnothing)\) on \(M_\iota\).
If every \(e\in I_p\) with \(\tau,(-p^e)\tau\in\Phi_a\) satisfies \(e\ge e_0\),
the geometric series gives \(C_p(\varnothing)\le 2fa/p^{e_0}\).
When \(p\equiv\pm1\bmod m\), take \(e_0=1\).
In the other case, a collision can occur only for inert \(p\), and then
\(e\) is an odd multiple of \(o_0\), the order of the place label of
\(-p\) modulo \(q_0\). If \(o_0>\lceil\log_p m\rceil\), one can take
\(p^{e_0}\ge pm\). In these cases \(mC_p(\varnothing)\le\ell_p\)
and \(M'_\iota=M_\iota\cap S\) suffices.

The remaining exceptional case is when \(p\) is inert,
\(p\not\equiv\pm1\bmod m\), and \(o_0\) is odd with
\(o_0\le\lceil\log_p m\rceil\).
Among the finite places of \(E\) induced by \(\Phi_a\) and their conjugates,
one can select one place from each conjugate pair to obtain a full
marking \(\Lambda\). If \(A\) is the set of infinite places inducing
places in \(\Lambda\), then \(C_p(A)=0\).
\leanref{exists_isFullMark}

Put \(r_p=\lceil\log_p m\rceil\) and \(\nu_p=r_pf(n_m+r_p)\).
For \(w\in\Lambda\), let \(a_w\) be \(a\) times the number of members of
\(\Phi_a\) inducing \(w\).
Let \(Y\) be the product over \(w\in\Lambda\) of the operators \(T_w\)
obtained by multiplying the Hecke operator of index \(r_p\) at \(w\)
by \(p^{a_w}\).
\emph{The marked estimate on neighboring lattices.}
Consider translating the kernel on the right at each \(w\in\Lambda\)
by an element of a Hecke double coset of index at most \(r_p\).
Up to integral factors on both sides, its component at \(w\)
is diagonal with \(p\) on a set of rows \(J_w\), \(|J_w|\le r_p\),
and \(1\) elsewhere. The conjugate component is the inverse.
Since the weight embedding reads the conjugate component, the row
factor has absolute value \(p^{\eta_{\tau,b}}\).
Here \(\eta_{\tau,b}\) is \(1\) if \(\tau\) induces \(w\in\Lambda\)
and \(b\in J_w\), is \(-1\) if \(\tau\) induces \(\overline w\)
and \(b\in J_w\), and is otherwise \(0\).
The determinant factor has exponent \(\sum_\tau e'_\tau\sum_b\eta_{\tau,b}\),
which vanishes by the equality of the sums of the last weight
entries above \(w\) and \(\overline w\).

For \eqref{eq:upper-jet}, take \(w_{\tau h b}=1\) if \(\tau\) induces
\(w\) and \(h=1,\ b\notin J_w\), or if \(\tau\) induces \(\overline w\)
and \(b\in J_w\). Take it to be \(0\) elsewhere.
The prime \(p\) is inert, and each relevant \(e\) is odd, so
\(\tau_e\) induces the conjugate of the prime induced by \(\tau\).
In each exempt pair, if \(b\in J_w\) then the conjugate side
receives a factor of \(p\). If \(b\notin J_w\), the marked column
of height \(1\) receives that factor. Thus the terms exempted from
\eqref{eq:upper-direction-pairing} also have the required divisibility.

The sum \(\eta_{\tau,b}+w_{\tau h b}\) for a box is zero at a place
inducing \(\overline w\). At a place inducing \(w\), a box in a
column of height \(1\) contributes \(1\). In any other column,
the row indices in a standard filling are distinct, so there are
at most \(r_p\) boxes with index in \(J_w\).
The sum over boxes is therefore at most
\(\sum_{w\in\Lambda}a_w+r_p\sum_{\tau\in A}r_\tau\).
Subtracting \(D_{p,j}(A)\) gives the upper bound on the loss
for the kernel on the neighboring lattices.
\leansrc{New/UpperBound/Divisibility/ValuationAtNonSplitPrimes}{Weight factors and powers of \(p\) for the neighboring lattices}.

\emph{Transferring the Hecke operators into the data.}
Let \(C_{\boldsymbol x,j}\) be the composite of the first \(j\) lifts
with the translation tuple \(\boldsymbol x\). For \(w\in\Lambda\),
write \(Q_w=z_w^2\) for the order of its residue field, with
\(z_w=p^{h_w}\) and \(2h_w\le f\). Denote the Hecke operator on
\(X_j\) of index \(i\) by \(T_{w,i}^{(j)}\).
Write \(t_{\rm rk}+1\) for the rank increment \(n_{j+1}-n_j\).
The character of the twist is unramified at this place, and the
local datum is the vacuum. For the Hecke polynomial
\[
 {\cal H}_{w,j}(U)=\sum_{i=0}^{n_j}(-U)^i Q_w^{\binom{i}{2}}T_{w,i}^{(j)}
\]
the identity for an unramified theta lift \(L_j\) is
\begin{equation}\label{eq:upper-hecke-polynomial}
 {\cal H}_{w,j+1}(U)L_j
    =R_j(U)L_j{\cal H}_{w,j}(z_w^{t_{\rm rk}+1}U).
\end{equation}
Here \(\chi_{w,j+1}\) is the value of the base change of the lifting
character at a uniformizer of \(w\), and
\(R_j(U)=\prod_{0\le k\le t_{\rm rk}}(1-Q_w^k z_w^{n_j}\chi_{w,j+1}U)\).
The equality of the sums of the character weights above \(w\) and
\(\overline w\) makes \(\chi_{w,j+1}\) a \(p\)-adic unit.
Hence \(R_j\in\mathcal O[U]\) and \(R_j(0)=1\).
\leansrc{Lib/Weil/Unramified/HeckeTransfer}{Hecke polynomial identity for the unramified theta lift}.

Multiply \eqref{eq:upper-hecke-polynomial} by \(R_j^{-1}\in\mathcal O[[U]]\)
and compare coefficients of \(U^i\).
Then \(z_w^{(t_{\rm rk}+1)i}Q_w^{\binom{i}{2}}L_j T_{w,i}^{(j)}\)
is an integral linear combination of \(T_{w,i'}^{(j+1)}L_j\)
for \(i'\le i\).
The prefactor is \(p^{h_w i(t_{\rm rk}+i)}\).
For \(i\le r_p\le t_{\rm rk}\), it divides \(p^{f r_p t_{\rm rk}}\),
since \(2h_w\le f\). Paying this power of \(p\) at each step thus
keeps the transfer integral and preserves the upper bound \(r_p\)
on the index. Iterate through \(j+1\) steps and multiply by the
remaining power of \(p\), using \((j+1)f r_p t_{\rm rk}\le\nu_p\).
This gives \(p\)-integral coefficients \(c_{w,i}\) with
\begin{equation}\label{eq:upper-hecke-transfer}
 \sum_{i\le r_p}c_{w,i}\,T_{w,i}^{(j+1)}C_{\boldsymbol x,j+1}
       = p^{\nu_p} C_{\boldsymbol x,j+1} T_{w,r_p}^{(0)} .
\end{equation}
The identity holds on forms spherical at \(w\), and its coefficients
are independent of \(\boldsymbol x\).
\leansrc{New/UpperBound/Divisibility/ValuationAtNonSplitPrimes}{Iterating the Hecke transfer through the lifts}.

Apply the operator on the left of \eqref{eq:upper-hecke-transfer}
to the datum and multiply by \(p^{a_w-\nu_p}\).
Do this at each of the distinct places \(w\).
Covariance moves the target Hecke action into the datum.
The datum is spherical at each \(w\), so the resulting operator
commutes with the target level and with the source action.
Writing \(L'_{\boldsymbol x,j}\) for the lift with the changed datum,
we have \(L'_{\boldsymbol x,j}C_{\boldsymbol x,j}
=C_{\boldsymbol x,j+1}Y\).
Induction gives \(C_{\boldsymbol x,m}Y^m\) for the composite of the
\(m\) changed lifts.

Each \(T_{w,i}^{(j+1)}\) is a finite sum of translations by one-sided
coset representatives in the double coset.
They all satisfy the conditions of the neighboring-lattice estimate.
The ultrametric triangle inequality preserves the bound in a
finite sum with the integral coefficients \(c_{w,i}\).
Multiplication by \(p^{a_w-\nu_p}\) adds a loss of \(\nu_p-a_w\)
for each \(w\). Hence the total loss in step \(j\) is at most
\[
  |\Lambda|\nu_p+r_p\sum_{\tau\in A}r_\tau-D_{p,j}(A).
\]
Commutation with the levels preserves joint invariance of the data
in the chain, so these losses can be summed over the steps.\leanref{gaugeLe_scaledHeckeOperator_pow}

The cost at the marked places satisfies
\[
 |\Lambda|\nu_p+r_p\sum_{\tau\in A}r_\tau
 \le \frac{\ell_p}{m}+|\Phi_a\cap A|\frac a{x_1}.
\]
\leansrc{New/UpperBound/Divisibility/MarkedPrimes}{Cost at marked places}.
At an exceptional prime \(v_p(\pi_p)=1\), so the terms involving
\(a/x_1\) cancel, and the loss per step is at most
\((2\ell_p-G_p)/m\).
Set \(j_p=\lceil 2mfa/p\rceil\).
On \(p^{j_p}M_\iota\), the valuation of the image is at least
\(G_p-\ell_p\), by the unmarked estimate and
\(mC_p(\varnothing)\le 2mfa/p\le j_p\).
Thus the bound \(G_p-2\ell_p\) holds on
\(Y^m M_\iota+p^{j_p}M_\iota\).

The normalized Hecke operator \(T_w\) preserves \(M_\iota\), and the
colength of \(T_wM_\iota+p^{j_p}M_\iota\) divides \(p^{c_w}\) for some
\[
  c_w\le\left(\frac{12a_w}{mp}+\frac{j_p}{m^2f}\right)\dim S_0 .
\]
\leansrc{New/UpperBound/Divisibility/MarkOperator}{Elementary divisors of the normalized Hecke operator}.
To see the numerical bound, the places of \(\Phi_a\) inducing
\(w\) or \(\overline w\) form one orbit under \(-p^{o_0}\).
Let its length be \(2u_w\) and put \(P_p=p^{o_0}\).
Order the embeddings of the orbit into the valuation ring, composing with
complex conjugation on the conjugate side. Successive embeddings are
related by the \(P_p\)-power map modulo \(p\). The even embeddings induce
\(w\), so \(a_w=au_w\).
The elementary-divisor theorem for a Frobenius cycle factors the Hecke operator on the lattice classes into a restriction
to residue classes and a pairing of matrices, and compares the
colengths of the images of these matrices.
For the coordinate model and numerical conditions
\[
 \begin{gathered}
  0<r_p\le k/2,\qquad a^4\le P_p^{2u_w},\qquad \theta_0=(m^2f)^{-1},\\
  (2u_w k\log(P_p^{2u_w})/\theta_0)^5 P_p^{3r_p}\le a
 \end{gathered}
\]
the bound per coordinate is \(3C_{\mathrm{div}}+\theta_0j_p\), where
\(C_{\mathrm{div}}=2u_wa(P_p^{-2r_p}+r_p^2/(kP_p))\).
\leansrc{New/UpperBound/ElementaryDivisors/ElementaryDivisorTheorem}{Elementary divisors for a Frobenius cycle}.
For the present Hecke operator, the hypotheses follow from its coordinate
model and the scale conditions.
\leansrc{New/UpperBound/Divisibility/MarkOperator}{Application of the elementary-divisor theorem}.
From \(m\le p^{r_p}\) and \(p\le P_p\) one gets
\(P_p^{-2r_p}\le1/(mp)\). From \(mr_p^2\le k\) and \(p\le P_p\) the
second term is at most \(1/(mp)\) as well.
Hence \(3C_{\mathrm{div}}\le12a_w/(mp)\).

Colength exponents add under composition of the operators.
First send a finite family into \(T_1M_\iota+p^{j_p}M_\iota\), then apply
a correcting matrix for \(T_2\) to the preimages under \(T_1\) of the
\(T_1M_\iota\) components. Since \(T_1\) preserves \(p^{j_p}M_\iota\),
the composite matrix sends the family into
\(T_1T_2M_\iota+p^{j_p}M_\iota\). Its determinant is the product of the
two determinants. Repeating this argument and using
\(\sum_w a_w\le af\), \(|\Lambda|\le f\), and \(j_p\le2mfa/p+1\), the
colength exponent of \(Y^mM_\iota+p^{j_p}M_\iota\) is bounded above by
\[
  m\sum_{w\in\Lambda}
      \left(\frac{12a_w}{mp}+\frac{j_p}{m^2f}\right)\dim S_0
    \le \frac{14fa+p/m}{p}\dim S_0
    \le \ell_p\dim S_0 .
\]
Intersecting with \(S\) preserves this bound because the original family
can be taken in \(S\) and its linear combinations then remain in \(S\).

Finally, for \(p=N\), take the submodule on which
\(\mathrm{Up}^{m/2}\phi\) is integral in the modified standard lattice.
Its colength exponent is at most \(mfb\dim S_0\).
\leansrc{New/UpperBound/Divisibility/ColengthAtLevel}{Colength at the level prime}.
The ramification condition expresses each glued pair of lifts as a
composite of \(\mathrm{Up}\) and direct lifts. On this submodule the
resulting valuation is at least \(-2mfb\).
\leansrc{New/UpperBound/Divisibility/ValuationAtLevel}{Valuation at the level prime}.
Intersecting with \(S\) preserves the colength bound.
Since \(G_N=0\), this proves \eqref{eq:upper-local} in all cases.

The local errors satisfy
\begin{equation}\label{eq:upper-error-sum}
  \sum_p \varepsilon_p\log p\le E_3.
\end{equation}
For \(p\in\mathcal U\), use \(\log p\le\log x_1\).
When \(p\equiv\pm1\bmod m\), write \(p=mc\pm1\) with \(1\le c\le p\)
and use \(m/p\le2/c\). Otherwise take \(c=p\).
For the term \(p=N\), use \(c=1\) and
\(2mfb\log N\le4fa(\log b)^2\log x_1\).
At most four terms correspond to a given \(c\), from
\(N,mc-1,c,mc+1\). Thus the sum is at most
\(16fa(\log b)^2\log x_1(1+\log x_1)\), which is bounded by \(E_3\)
by the scale conditions. \leanref{sum_allowedLoss_mul_log_le}

\subsection{Archimedean norms and normalization}

For step \(j\) of the upper chain, let \(l_{j,\tau,i}\), \(i<n_j\),
be the row lengths used at \(\tau\) and put
\(U_j=\sum_\tau\sum_{i<n_j}s_\infty(l_{j,\tau,i})\).
Recall that \(s_\infty(q)=\log(q!)+q\log(M/\Omega)\), where \(\Omega\)
is the period. The row sum \(U_j\) gives the factorial contribution
in the norm of the kernel.
Use the Gram norm on the Weyl modules and test the Schur matrix
against row vectors of dual norm at most \(1\).
Let \(\rho_{\tau,i}\ge0\) be the squared singular values of
\(Z_\tau\), in decreasing order. The squared operator norm of the Schur matrix is
at most \(\prod_{\tau,i}\rho_{\tau,i}^{l_{j,\tau,i}}\).
Indeed, the unitary factors preserve the Gram norm, and for the
diagonal factor the highest weight gives the exponents \(l_{j,\tau,i}\).
The sum of the squared singular values is the sum of squared
absolute values of the entries of \(Z_\tau\).
\leansrc{Lib/GLn/SchurModule/SchurNorm}{Singular-value bound for the Schur matrix}.

Put \(\lambda=a+1\), \(u=\log(M/\Omega)\), and \(v=u+1/\lambda\).
Taking the term of degree \(l\) in the exponential series gives
\(\rho^l\le l!\exp(lv+\rho e^{-v})\).
If \(\sum_{\tau,i}\rho_{\tau,i}\le e^u(2\lambda+2)R\), then
\(e^{-v}\sum_{\tau,i}\rho_{\tau,i}\le2\lambda R\), since
\(e^{1/\lambda}\ge1+1/\lambda\).
Thus the Schur polynomial satisfies a growth bound with factor
\[
 \exp\!\left(\frac{U_j}{2}
       +\frac1{2\lambda}\sum_{\tau,i}l_{j,\tau,i}+\lambda R\right).
\]
In the conjugated Gaussian integral, the scalar multiplying the
variables has squared modulus \(M/\Omega\).

The conjugated derivative is a Gaussian integral of the Schur
polynomial, with its places permuted, against a conjugate
translated vacuum. With \(s_j=f n_j n_{j+1}\) complex variables \(y\),
put \(R=\sum|y|^2/(2\lambda+2)\).
The factor \(\exp(\lambda R)\) from polynomial growth times
\(\exp(-\sum|y|^2/2)\) from the vacuum is \(\exp(-R)\).
The measure is Lebesgue measure divided by \(\pi\) in each complex
coordinate, so the integral of this Gaussian is \((2\lambda+2)^{s_j}\).
It suffices to use the integral representation for conjugates fixing
\(K\), since composing with complex conjugation preserves the
absolute values for the other conjugates.
\leansrc{New/UpperBound/Kernels/JetConjugatesAndBounds}{Gaussian integral and the row factorials}.

The logarithmic constant from the bound on the conjugate of
the integral vacuum is \(\mathrm{const}(P)(q_1 n_j n_{j+1}+1)^2\).
The sum of this constant, \(s_j\log(2\lambda+2)\), and
\((2\lambda)^{-1}\sum_{\tau,i}l_{j,\tau,i}\) is at most \(x^{50}\)
by the bound \(2a\) on the row lengths and the scale conditions.
Thus a single lift increases the norm by at most
\(\exp(x^{50}+U_j/2)\) for every conjugate.\leanref{sqrt_weylNormSq_jetMap_le}

The principal contributions to the row lengths are \(a\,1_{\Phi_a}(\tau)\)
in the first row and \(\kappa(\tau)\) in each of the next \(k-1\) rows.
Each additional box contributes at most \(25\log b\). Indeed, all the
row lengths are at most \(2a\) and adding a box increases the bound by at
most \(\log(2a+1)+\log(M/\Omega)\le25\log b\).
At each step the total number of additional boxes is at most
\(f(b+a/x_1+n_m r^{16})\).
The scale conditions give
\begin{equation}\label{eq:upper-rows}
 \sum_{0\le j<m}U_j
 \le m\left\{|\Phi_a|s_\infty(a)+k\sum_\tau s_\infty(\kappa(\tau))\right\}
         +\frac{E_3}{2}.
\end{equation}
\leanref{sum_rowSizeSum_le}

In the expansion of the glued tensor for the trivial twist, each step
has at most \(\exp(x^{30})\) classes, and the sum of the absolute values
of the conjugated translation coefficients is at most \(\exp(x^{60})\).
If all values of \(\phi\) have squared norm at most \(\exp(B)\) after
conjugation, summing the bounds for products of the \(m\) kernels
bounds the squared norms before the final restriction by
\[
  \exp\!\left(B+\sum_j U_j
             +2(mx^{30}+x^{60}+mx^{50})\right).
\]
The restriction at the last step uses a branching matrix with
integer entries, which multiplies the
squared norm by at most \(\exp(x^{70})\).
Write \(u_\tau\) for the difference of the sums of the weights of
\(X_m,Y_m\) at \(\tau\). The restriction to \(Y_m\subset X_m\)
also multiplies by \(\prod_\tau\tau(\overline{\pi_N})^{-u_\tau}\).
Since \(\pi_N\overline{\pi_N}=N\), this factor has squared absolute value
\(N^{-\sum_\tau u_\tau}\) at every conjugate. The exponent therefore
decreases by at least \((|\Phi_a|a-fmr^8)\log N\).
\leansrc{New/UpperBound/Divisibility/ChainValuation}{The final restriction}.

The remainder is
\(\mathcal R=\sum_{\tau\in\Phi_a}(s_{\mathrm{fin}}(a,\tau)-s_{\mathrm{main}})\),
where \(s_{\mathrm{main}}\) is the main term for \(s_{\mathrm{fin}}(a,\tau)\).
Writing \(\overline{\log\kappa}=f^{-1}\sum_\tau\log\kappa(\tau)\) and
\(D_E=\operatorname{rd}(E)\), the condition on the level is
\[
 \log N\ge m\!\left(\overline{\log\kappa}
                 +\log \frac{M D_E}{2\pi}\right)
          -\log\frac{M D_E}{2\pi}-\tfrac{15}2\log b.
\]
The period formula
\(\log(M/\Omega)=1+c_1+\log(M D_E/(2\pi))\),
the identity \(s_{\mathrm{main}}=a(\log a-\overline{\log\kappa}+c_1)\),
and Stirling's upper bound
\(\log(a!)\le a\log a-a+1+\frac12\log a\) give
\[
 s_\infty(a)-s_{\mathrm{main}}
 \le a\!\left(\overline{\log\kappa}+\log \frac{M D_E}{2\pi}\right)
        +1+\tfrac12\log a .
\]
Consequently \(m(s_\infty(a)-s_{\mathrm{main}})-a\log N\le32a\log b\).
This uses \(m(1+\frac12\log a)\le a\log b\) and
\(\log(M D_E/(2\pi))\le23\log b\).
Substituting \eqref{eq:upper-gain-sum} gives
\[
 H_0-m\mathcal R+2\sum_p G_p\log p
 =m|\Phi_a|s_{\mathrm{main}}+mk\sum_\tau s_\infty(\kappa(\tau)).
\]
Use this identity and \eqref{eq:upper-rows} and absorb
\(32fa\log b+x^{70}+fmr^8\log N+
  2(mx^{30}+x^{60}+mx^{50})\) into \(E_3/2\).
All conjugated values of \(B_{1,\Psi}\phi\) then have squared norms at most
\begin{equation}\label{eq:upper-arch}
  \exp\!\left(B+H_0-m\mathcal R+2\sum_p G_p\log p+E_3\right).
\end{equation}
\leanref{weylNormSq_restrictedChain_le}

\paragraph{A common rational normalization}
Put \(g_p=\lfloor G_p-\varepsilon_p\rfloor\) and set
\[
 C_B=\prod_{p\in\mathcal U\cup\{N\}}p^{-g_p},\qquad
 \mathfrak p_g\phi
   = C_B|\overline G|e_0\rho(g)c_{\mathrm{ss}}^{-1}\mathcal P\phi .
\]
The factor \(C_B\) is a positive rational number.
By \eqref{eq:upper-local}, \(C_B B_{\xi,\Psi_\xi}\phi\) is integral on
\(M'_\iota\), and \eqref{eq:upper-seesaw} and \eqref{eq:upper-project}
show that \(\Theta_\xi\mathfrak p_g\phi\) is integral as well.
The same \(M'_\iota\) works for every allowed twist.

The rounding costs at most \(\sum_{p\in\mathcal U\cup\{N\}}\log p\le b\).
Using \eqref{eq:upper-error-sum} we get
\begin{equation}\label{eq:upper-cb}
 \log C_B\le-\sum_p G_p\log p+E_3+b .
\end{equation}
Multiplication of the squared norm by \(C_B^2\) combines with the
finite-place gains in the exponent as follows.
\[
 2\log C_B+2\sum_p G_p\log p
    =2\sum_{p\in\mathcal U\cup\{N\}}(G_p-g_p)\log p
    \le 2(E_3+b).
\]
Thus the principal finite-place term in \eqref{eq:upper-arch} cancels,
leaving this error and the projection loss \(x^{99}\).

\begin{proof}[Proof of Theorem~\ref{thm:upper-bound}]
The space \(S_0\) is spanned by forms integral at all finite places
whose values, after any conjugation, have squared norm at most
\(\exp(E_3)/4^f\). To obtain these forms, use that the stabilizer of
each lattice class is trivial and take forms supported on one class
at a time. For a class \(c\), choose short bases for its lattice and
for the lattices entering the determinant part of the weight.
The evaluation vectors of the Weyl module at roots of unity span the
module, and applying the chosen basis matrices preserves their span.
The bases are locally integral and the roots of unity have absolute
value \(1\), so the resulting forms are integral at all finite places.
The archimedean constant for the short bases is \(C=6(q_1d)(fk)^2\).
At all places together, columns of height less than \(k\) contain at
most \(fa\) boxes. The squared-norm bound is then \(C^{6(fk+fa)}\).
Admissibility gives \(4^f C^{6(fk+fa)}\le\exp(E_3)\).
\leanref{span_isShort_eq_top}

Let \(z_w\) be the generator of the scalar action at a specified place
\(w\) above \(N\). The image of \(\prod_w(1-z_w)\) is \(S\).
Scalar translation preserves integrality, and each difference
\(1-z_w\) increases the squared-norm bound by a factor of at most \(4\).
There are \(f\) specified places, so the images of the preceding forms
span \(S\) and satisfy the bound with \(B=E_3\).
Combine the exponent in \eqref{eq:upper-arch} with the projection
loss \(x^{99}\) and twice the bound on \(\log C_B\) in
\eqref{eq:upper-cb} to obtain
\[
 \|\Theta_1\mathfrak p_g\phi\|_\sigma^2
 \le \exp\bigl(H_0-m\mathcal R+x^{99}+4E_3+2b\bigr).
\]
The admissibility inequality \(x^{99}+4E_3+2b\le E_4\) proves (1).

For (2), first fix \(g_i=g\).
For each \(\iota\), apply the correcting matrix from
\eqref{eq:upper-local} to \((\phi_i)_i\) to obtain a family in \(M'_\iota\).
The same matrix sends \((\mathfrak p_g\phi_i)_i\) to a family integral
for every allowed twist. The valuation of its determinant is at most
\(\varepsilon_p\dim S_0\).
Imposing ramification loses at most \(\frac{f}{N-1}\dim S_0\)
dimensions, and \(\frac{f}{N-1}\le1/2\), so \(\dim S_0\le2\dim S\).
By \eqref{eq:upper-error-sum}, the logarithmic index for a fixed \(g\)
is at most \(2(\dim S)E_3\).
For a general family, group its members according to \(g_i\) and take
the block diagonal matrix with the corresponding correcting matrices
as blocks. Its determinant is the product of the determinants of the
blocks. Adding their logarithmic index bounds gives
\(2|\mathcal T|(\dim S)E_3\).
\end{proof}
